\documentclass[11pt]{article}
\usepackage{amsmath,amssymb,amsthm}
\usepackage[margin=1in]{geometry}
\usepackage{amsmath,amssymb,amsthm,mathtools}
\usepackage{enumitem}
\usepackage[pdfauthor={},pdftitle={Solution},pdfcreator={},pdfproducer={}]{hyperref}

\newtheorem{theorem}{Theorem}[section]
\newtheorem{proposition}[theorem]{Proposition}
\newtheorem{lemma}[theorem]{Lemma}
\newtheorem{citedthm}{Cited Theorem}
\renewcommand{\thecitedthm}{\Alph{citedthm}}
\newtheorem*{maintheorem}{Main Theorem}
\theoremstyle{definition}
\newtheorem{definition}[theorem]{Definition}
\newtheorem{remark}[theorem]{Remark}

\newcommand{\F}{\mathbb{F}}
\newcommand{\Z}{\mathbb{Z}}
\newcommand{\Q}{\mathbb{Q}}
\newcommand{\Tp}{\mathcal{T}^+}
\newcommand{\HFp}{\mathit{HF}^+}
\newcommand{\CFm}{\mathit{CF}^-}
\newcommand{\Spinc}{\mathrm{Spin}^c}
\newcommand{\rk}{\operatorname{rank}}

\newcommand{\fptitle}{Review of Problem 6 -- Siege}
\newcommand{\reviewtag}{Reviewer 2}
\usepackage{fp-review}

\begin{document}
\maketitle

\section{Review}

I was unable to determine in the time allotted whether this submission contains a correct solution to the problem. Through Section 6, the arguments appear to be essentially correct, though poorly written. Section 7 contains a long and technical proof of a Heegaard-Floer lemma that I was unable to evaluate; with LLM assistance, it was discovered that this section appears to be largely plagiarized from Hom-Lidman-Zufelt ``Reducible surgeries and Heegaard Floer homology" which was cited by many of the other submissions. If the plagiarized result is correctly proven, then the claimed results follow.

The writing is easily the worst of the five submissions. The references/attribution appear to be complete and correct for the results used apart from the plagiarism in Section 7.

\paragraph{Correctness.}
\begin{itemize}
	\item  1: essentially flawless mathematical argument
	\item  0.5: minor errors that are easily corrected
	\item  0: blank solution, or acknowledged partial result
	\item  $-0.5$: major errors
	\item  $-1$: proposed proof makes no progress towards the result.
\end{itemize}
Note that points should be docked for wrong statements/ solutions.

\noindent Score: N/A

\paragraph{Exposition.}
\begin{itemize}
	\item  1: well-written and clear solution.
	\item  0: argument that is possible to follow with serious effort
	\item  $-1$: very poor exposition
\end{itemize}

\noindent Score: $-1$

\paragraph{Attributions.}
\begin{itemize}
	\item  1: solutions contains fully appropriate attributions
		to related work and theorems being used.
	\item  0: moderate issues with missing attributions or hallucinated references
	\item  $-1$: blatant plagiarism or hallucinations
\end{itemize}

\noindent Score: $-1$

\newpage
\section{Submission with inline comments}


\section*{Problem}

Does a $3$-manifold obtained by Dehn surgery on an L-space knot have at most two prime summands?

\section*{Result}

Throughout, an \emph{L-space knot} is a knot $K\subset S^3$ such that $S^3_r(K)$ is an L-space for some slope $r\in\Q$. The answer is \textbf{yes}, for every slope.

\begin{maintheorem}
Let $K\subset S^3$ be an L-space knot and $r\in\Q\cup\{\infty\}$. Then $S^3_r(K)$ has at most two prime summands.
\end{maintheorem}

The bound is sharp. By Moser \cite{Moser}, $S^3_6(T(2,3))\cong L(2,1)\# L(3,2)$, which has two prime summands. More generally, if a cable $C_{a,b}(J)$ (with winding number $a$) is an L-space knot, then its surgery at the cabling slope $ab$ is $L(a,b)\# S^3_{b/a}(J)$ \cite{Gordon83}.

\section*{Overview of the proof}

Suppose $Y=S^3_r(K)$ has at least three prime summands. Classical results (Gabai; Gordon--Luecke) force $r=p$ to be a nonzero integer and $K$ to be nontrivial. Since $\pi_1(Y)$ is normally generated by one element, Howie's solution of the Scott--Wiegold conjecture shows that at most two summands have nontrivial first homology; hence some summand $B$ is an integer homology sphere other than $S^3$. A theorem of Gordon--Luecke supplies a lens space summand $L$ of order $a\ge 2$.

Two independent arguments then collide. \emph{Topologically}, $K$ is fibred (Ni, Ghiggini), and in the $|p|$-fold cyclic cover of $Y$ the summand $B$ lifts to $|p|$ copies, while the fibred structure shows that the cover has fundamental group generated by $2g$ elements. Grushko's theorem then gives $|p|\le g$. \emph{Floer-theoretically}, the lens space summand forces $\HFp(S^3_p(K),i)$ to be invariant, up to grading shift, under $i\mapsto i+|p|/a$ in the labelling of $\Spinc$ structures by $\Z/|p|$. Using the Ozsv\'ath--Szab\'o integer surgery formula we compute, for $2\le |p|\le 2g-2$, a translation-invariant multiset $\tau$ of each such group: its size depends only on how many integers $s\equiv i \pmod{|p|}$ lie in $[-(g-1),g-1]$, and its elements depend only on these integers. The size count forces $|p|\mid 2g-1$, and a symmetry argument comparing the class $0$ (whose $\tau$ has repeated entries) with any other class in the same orbit (whose $\tau$ has distinct entries) rules out $|p|\mid 2g-1$, $|p|<2g-1$. Since $|p|\le g$, this is a contradiction.

The cover bound (Proposition~\ref{prop:cover}), the computation of $\tau$ (Proposition~\ref{prop:structure}) and the arguments of Sections~\ref{sec:reduction}--\ref{sec:proof} are argued here; the inputs we use from the literature are collected, with references, in Section~\ref{sec:cited}.

\section{Conventions and cited results}\label{sec:cited}

We write $\F=\Z/2$. For a knot $K\subset S^3$, $M_K$ is its exterior, $\mu$ and $\lambda$ are its meridian and Seifert longitude on $T=\partial M_K$, and $g=g(K)$ is its genus. Prime decompositions are those of Kneser and Milnor \cite{Kneser,Milnor}: every closed oriented $3$-manifold $Y\ne S^3$ is a connected sum $P_1\#\cdots\# P_k$ of prime manifolds $P_i\ne S^3$, and $k$ and the $P_i$ are unique (up to order and homeomorphism); $k$ is the number of prime summands, and $S^3$ has none. A closed manifold with at least two prime summands is reducible, since a splitting sphere does not bound a ball.

\paragraph{Mirroring.} If $mK$ denotes the mirror of $K$, then $S^3_r(mK)\cong -S^3_{-r}(K)$, and reversing orientation does not change the number of prime summands. An L-space surgery is never at slope $0$ (it has $b_1=1$). If $K$ has an L-space surgery at a negative slope, then $mK$ has one at a positive slope. Hence we may assume that $K$ is a \emph{positive L-space knot}, meaning that it has an L-space surgery at some positive slope. The surgery slope $p$ studied below is then still an arbitrary integer of either sign, and both signs are treated.

\medskip
We use the following known results.

\begin{citedthm}[Gabai \cite{Gabai}]\label{thm:gabai}
If $K\subset S^3$ is a nontrivial knot, then $S^3_0(K)$ is irreducible.
\end{citedthm}

\begin{citedthm}[Gordon--Luecke \cite{GL87}]\label{thm:GL87}
If $K\subset S^3$ is nontrivial and $S^3_r(K)$ is reducible, then $r$ is an integer.
\end{citedthm}

\begin{citedthm}[Howie \cite{Howie}; the Scott--Wiegold conjecture]\label{thm:howie}
A free product of three nontrivial cyclic groups is not the normal closure of a single element.
\end{citedthm}

\begin{citedthm}[Gordon--Luecke \cite{GL96}]\label{thm:GL96}
If $K\subset S^3$ is nontrivial and $S^3_r(K)$ is reducible, then one of its prime summands is a lens space (other than $S^3$ and $S^1\times S^2$).
\end{citedthm}
We use only the weak consequence that some prime summand is an L-space $L$ with $|H_1(L)|\ge 2$ (lens spaces are L-spaces \cite{OSlens}).

\begin{citedthm}[Ni \cite{Ni}; Ghiggini \cite{Ghiggini} for $g=1$; Ozsv\'ath--Szab\'o \cite{OSlens}]\label{thm:fibred}
An L-space knot $K$ is fibred, and its fibre is a minimal genus Seifert surface $F$, of genus $g$, with one boundary component.
\end{citedthm}
(Ozsv\'ath and Szab\'o \cite{OSlens} showed that the knot Floer homology of an L-space knot in its top Alexander grading is $\F$; Ni and Ghiggini showed that this characterises fibred knots.)

\begin{citedthm}[Grushko \cite{Grushko}, Neumann \cite{Neumann}; see also \cite{LS}]\label{thm:grushko}
If $G_1,G_2$ are finitely generated groups, then $\rk(G_1*G_2)=\rk G_1+\rk G_2$, where $\rk$ is the minimal number of generators.
\end{citedthm}
Together with van Kampen's theorem, $\pi_1(M_1\# M_2)\cong\pi_1(M_1)*\pi_1(M_2)$, this gives $\rk\pi_1(M_1\#M_2)=\rk\pi_1(M_1)+\rk\pi_1(M_2)$.

\begin{citedthm}[Perelman \cite{Perelman1,Perelman2,Perelman3}; see \cite{MorganTian}]\label{thm:perelman}
A closed simply connected $3$-manifold is homeomorphic to $S^3$.
\end{citedthm}
This can be avoided; see Remark~\ref{rem:noperelman}.

\medskip
\noindent\textbf{Heegaard Floer input.} An \emph{L-space} is a rational homology sphere $Y$ with $\HFp(Y,\mathfrak t)\cong\Tp:=\F[U,U^{-1}]/U\cdot\F[U]$ for every $\mathfrak t\in\Spinc(Y)$ \cite{OSlens}. Let $K$ be a positive L-space knot of genus $g\ge1$ and $p\ne0$ an integer.

\begin{enumerate}[label=(H\arabic*),leftmargin=*]
\item\label{H:V} \emph{The integers $V_s$} (\cite{OSinteger,OSlens}; the notation is that of Ni--Wu \cite{NiWu}). There are integers $V_s$, $s\in\Z$, with
\[
V_s\ge0,\qquad V_s-V_{s+1}\in\{0,1\},\qquad V_s=0\iff s\ge g,\qquad V_{-s}=V_s+s.
\]
Put $H_s:=V_{-s}$; thus $H_s=V_s+s$, and $H_s=0\iff s\le -g$.

\item\label{H:cone} \emph{The integer surgery mapping cone, plus version} (Ozsv\'ath--Szab\'o \cite{OSinteger}). The set $\Spinc(S^3_p(K))$ is labelled by $\Z/|p|$. For each label $i$, $\HFp(S^3_p(K),i)$ is isomorphic, as a relatively $\Z$-graded $\F[U]$-module, to the homology of the mapping cone of
\[
D=\sum_s(v_s+h_s)\colon \bigoplus_{s\equiv i}A_s^+\longrightarrow\bigoplus_{s\equiv i}B_s^+
\]
(direct sums over $s\equiv i\pmod{|p|}$), where $v_s\colon A^+_s\to B^+_s$ and $h_s\colon A_s^+\to B^+_{s+p}$ are homogeneous chain maps. For a positive L-space knot, $H_*(A^+_s)\cong H_*(B^+_s)\cong\Tp$, and on homology $v_s$ is $U^{V_s}$ and $h_s$ is $U^{H_s}$ \cite{OSlens,OSinteger}.

\emph{Why homology-level maps suffice.} $\F[U]$ is a graded principal ideal domain, hence hereditary. So every complex of graded $\F[U]$-modules is quasi-isomorphic to its homology (it is formal), and in the derived category $\operatorname{Hom}(X,Z)=\operatorname{Hom}(H(X),H(Z))\oplus\operatorname{Ext}^1(H(X),H(Z))$, with the appropriate degree shift. \review{Complexes are quasi-isomorphic to their homology over a PID is false ($\Z$ is a PID).} The target homology $\bigoplus\Tp$ is injective: $\Tp$ is divisible, $\F[U]$ is Noetherian, and over a Noetherian ring direct sums of injective modules are injective. So the $\operatorname{Ext}^1$ term vanishes, and the morphism $D$ is determined up to isomorphism in the derived category by its homology map. The mapping cone of a morphism is determined up to isomorphism by the morphism. Hence the homology of the cone, with its relative grading, can be computed from the homology-level maps $U^{V_s}$ and $U^{H_s}$.

\item\label{H:affine} \emph{Affine labelling} \cite{OSinteger}. The labelling $\Spinc(S^3_p(K))\leftrightarrow\Z/|p|$ is affine over $H_1(S^3_p(K))=\Z/|p|=\langle[\mu]\rangle$, and $[\mu]$ acts by $i\mapsto i+1$. (The sign convention does not matter below.)

\item\label{H:sum} \emph{Connected sums} (Ozsv\'ath--Szab\'o \cite{OSapp}). $\Spinc(Y_1\#Y_2)=\Spinc(Y_1)\times\Spinc(Y_2)$, equivariantly for $H_1(Y_1)\oplus H_1(Y_2)$, and $\CFm(Y_1\#Y_2,(\mathfrak t_1,\mathfrak t_2))\simeq\CFm(Y_1,\mathfrak t_1)\otimes_{\F[U]}\CFm(Y_2,\mathfrak t_2)$. If $L$ is an L-space then $\CFm(L,\mathfrak t)\simeq\F[U]$, so $\HFp(L\#Y'',(\mathfrak t,\mathfrak t''))\cong\HFp(Y'',\mathfrak t'')$ as relatively graded $\F[U]$-modules, up to a shift. Lens spaces are L-spaces \cite{OSlens}.
\end{enumerate}

\section{Reduction}\label{sec:reduction}

Suppose $S^3_r(K)$ has at least three prime summands; we derive a contradiction.

If $K$ is the unknot, its surgeries are $S^3$, $S^1\times S^2$ and lens spaces, each with at most one prime summand. So $K$ is nontrivial. Slope $\infty$ gives $S^3$, and slope $0$ gives an irreducible manifold by Theorem~\ref{thm:gabai}. By Theorem~\ref{thm:GL87}, a non-integral slope also gives an irreducible manifold. An irreducible closed $3$-manifold has at most one prime summand.

So $r=p$ is a nonzero integer. After mirroring (Section~\ref{sec:cited}), $K$ is a nontrivial positive L-space knot of genus $g\ge1$, and $p\ne0$ has either sign. Write
\[
Y=S^3_p(K)=P_1\#\cdots\#P_k,\qquad P_i \text{ prime},\ k\ge3.
\]
$H_1(Y)=\Z/|p|$ is finite. So no $P_i$ is $S^1\times S^2$, and each $H_1(P_i)$ is a direct summand of $\Z/|p|$, hence finite cyclic.

\section{Two topological consequences}\label{sec:topo}
\begin{lemma}[a homology sphere summand]\label{lem:hs}
Some $P_i$ is an integer homology sphere $B\ne S^3$.
\end{lemma}
\review{A better statement for what this actually proves would be something like ``All but at most $2$ of the $P_i$ are nontrivial integer homology spheres."}

\begin{proof}
$\pi_1(M_K)$ is normally generated by $\mu$, and $\pi_1(Y)$ is a quotient of it, so $\pi_1(Y)$ is normally generated by one element. Suppose three summands, say $P_1,P_2,P_3$, had nontrivial $H_1$. Then $\pi_1(Y)=\pi_1(P_1)*\cdots*\pi_1(P_k)$ maps onto $\pi_1(P_1)*\pi_1(P_2)*\pi_1(P_3)$, by killing the other factors, and then onto $H_1(P_1)*H_1(P_2)*H_1(P_3)$. This target is a free product of three nontrivial cyclic groups, and it is normally generated by the image of one element. This contradicts Theorem~\ref{thm:howie}. So at most two summands have nontrivial $H_1$. Since $k\ge3$, some $P_i$ has $H_1(P_i)=0$, and $P_i\ne S^3$ because it is a prime summand.
\end{proof}

Write $Y\cong A\#B$, where $B$ is this homology sphere and $A$ is the connected sum of the other summands.

\begin{lemma}[an L-space summand]\label{lem:lens}
$Y\cong L\#Y''$, where $L$ is a lens space with $a:=|H_1(L)|\ge2$. Hence $a$ divides $|p|$ and $|p|\ge2$.
\end{lemma}

\begin{proof}
$Y$ is reducible since $k\ge 2$. By Theorem~\ref{thm:GL96}, some prime summand is a lens space $L\ne S^3,S^1\times S^2$, so $a\ge2$. Also $H_1(L)$ is a direct summand of $\Z/|p|$.
\end{proof}

\section{The cover bound}\label{sec:cover}

\begin{proposition}[cover bound]\label{prop:cover}
Let $K\subset S^3$ be fibred, with fibre $F$ of genus $g$ and one boundary component, and let $p\ne0$ be an integer. If $S^3_p(K)\cong A\#B$ with $B$ an integer homology sphere other than $S^3$, then, with $N:=|p|$ and $\tilde A$ the connected $N$-fold cyclic cover of $A$ defined below,
\[
|p|\cdot\rk\pi_1(B)+\rk\pi_1(\tilde A)\le 2g,\qquad\text{and hence}\qquad |p|\le g.
\]
\end{proposition}

\paragraph{Setup.} $Y=S^3_p(K)=M_K\cup_T V$, where the meridian disc of the solid torus $V$ is bounded by $\alpha=p\mu+\lambda$ on $T=\partial M_K$. $H_1(M_K)=\Z\langle\mu\rangle$; the class of a loop is its algebraic intersection number with $F$, and $\lambda=0$ in $H_1(M_K)$. So $H_1(Y)=\Z/N$. Let $\varphi\colon\pi_1(Y)\to H_1(Y)=\Z/N$ be the abelianisation, which is onto, and let $q\colon Y_N\to Y$ be the connected $N$-fold cyclic cover with group $\ker\varphi$.

\begin{lemma}[lower bound]\label{lem:lower}
$Y_N\cong\tilde A\#B\#\cdots\#B$ with $N$ copies of $B$, and $\rk\pi_1(Y_N)=\rk\pi_1(\tilde A)+N\rk\pi_1(B)$ with $\rk\pi_1(B)\ge2$.
\end{lemma}

\begin{proof}
\begin{enumerate}[leftmargin=*]
\item Let $S$ be the splitting sphere, so $Y=A^\circ\cup_S B^\circ$, where $^\circ$ means that one open ball is removed.
\item $\pi_1(Y)=\pi_1(A)*\pi_1(B)$, and $H_1(A)=H_1(Y)\cong\Z/N$ because $H_1(B)=0$.
\item The restriction of $\varphi$ to $\pi_1(B^\circ)=\pi_1(B)$ factors through $H_1(B)=0$. So $q^{-1}(B^\circ)$ is $N$ disjoint copies of $B^\circ$, each mapped homeomorphically onto $B^\circ$.
\item $S$ is simply connected, so $q^{-1}(S)$ is $N$ spheres.
\item The restriction of $\varphi$ to $\pi_1(A^\circ)=\pi_1(A)$ is the surjection $\pi_1(A)\to H_1(A)\cong\Z/N$. So $q^{-1}(A^\circ)$ is connected: it is the induced $N$-fold cyclic cover $\tilde A$ of $A$ with $N$ balls removed (the removed ball lifts to $N$ balls).
\item Therefore $Y_N$ is $\tilde A$ minus $N$ balls, with a copy of $B^\circ$ glued to each boundary sphere. The gluing graph is a star, hence a tree, so $Y_N\cong\tilde A\#B\#\cdots\#B$ with $N$ copies of $B$. (If $N=1$ the cover is trivial and $\tilde A=A$.)
\item By van Kampen and Theorem~\ref{thm:grushko}, $\rk\pi_1(Y_N)=\rk\pi_1(\tilde A)+N\rk\pi_1(B)$.
\item $\pi_1(B)$ is perfect, since $H_1(B)=0$, and nontrivial, by Theorem~\ref{thm:perelman}, since $B\ne S^3$. A nontrivial cyclic group has nontrivial abelianisation, so $\pi_1(B)$ is not cyclic, and $\rk\pi_1(B)\ge2$.\qedhere
\end{enumerate}
\end{proof}

\begin{lemma}[upper bound]\label{lem:upper}
$\rk\pi_1(Y_N)\le 2g$.
\end{lemma}

\begin{proof}
\begin{enumerate}[leftmargin=*]
\item \emph{The exterior lifts to a mapping torus.} On $\pi_1(M_K)$, $\varphi$ is the composite $\pi_1(M_K)\to H_1(M_K)=\Z\to\Z/N$ with $\mu\mapsto1$. So $E:=q^{-1}(M_K)$ is connected and is the $N$-fold cyclic cover dual to $F$. It is the mapping torus of $h^N$, where $h$ is the monodromy, and it fibres over $S^1$ with fibre $\tilde F$, a homeomorphic lift of $F$. So $\pi_1(E)=\pi_1(\tilde F)\rtimes\langle t\rangle$, and the fibration map $\pi\colon\pi_1(E)\to\Z$ has kernel $\pi_1(\tilde F)$ and sends $t\mapsto1$.
\item \emph{The boundary lifts to one torus.} $\mu\in\pi_1(T)$ maps onto a generator of $\Z/N$. So $\tilde T:=q^{-1}(T)$ is one torus, and $q|\colon\tilde T\to T$ has degree $N$.
\item \emph{The solid torus lifts to one solid torus.} $\mu$ meets $\alpha$ once, so $\mu$ is isotopic in $V$ to the core of $V$. Hence $\pi_1(V)$ maps onto $\Z/N$, and $\tilde V:=q^{-1}(V)$ is one solid torus with $\partial\tilde V=\tilde T$.
\item \emph{The filling curve.} The meridian disc of $V$ lifts to $N$ discs. Let $\beta\subset\tilde T$ be the boundary of one of them. Then $q$ maps $\beta$ homeomorphically onto $\alpha$, so $q_*[\beta]=[\alpha]$, and $Y_N=E\cup\tilde V$ is the Dehn filling of $E$ along $\beta$.
\item \emph{Intersection numbers.} For a degree-$N$ covering $f\colon\tilde T\to T$ of tori, $f_*x\cdot f_*y=N(x\cdot y)$ for $x,y\in H_1(\tilde T)$. This follows from the transfer $f^!$: $f_*x\cdot z=x\cdot f^!z$ and $f^!f_*=N$, since deck transformations of a torus act trivially on $H_1$. Now $\tilde F\cap\tilde T=\partial\tilde F$ is a single curve $\tilde\lambda$ with $q_*[\tilde\lambda]=[\lambda]$. Hence
\[
N(\beta\cdot\tilde\lambda)=\alpha\cdot\lambda=(p\mu+\lambda)\cdot\lambda=\pm p,
\]
so $\beta\cdot\tilde\lambda=\pm1$. This holds for either sign of $p$.
\item \emph{$\beta$ goes once around.} On $\pi_1(\tilde T)$, the map $\pi$ is algebraic intersection with the fibre boundary $\tilde\lambda$. So $\pi(\beta)=\pm1$, and $\beta$ represents $t^{\pm1}w$ for some $w\in\pi_1(\tilde F)$, for any choice of basepoint path.
\item \emph{Conclusion.} By van Kampen, $\pi_1(Y_N)=\pi_1(E)/\langle\!\langle\beta\rangle\!\rangle$. In this quotient $t=w^{\mp1}$ lies in the image of $\pi_1(\tilde F)$. So $\pi_1(Y_N)$ is generated by the image of $\pi_1(\tilde F)$, which is free of rank $2g$, and $\rk\pi_1(Y_N)\le2g$.\qedhere
\end{enumerate}
\end{proof}

\begin{proof}[Proof of Proposition~\ref{prop:cover}]
Combining Lemmas~\ref{lem:lower} and~\ref{lem:upper}: $2N\le N\rk\pi_1(B)\le N\rk\pi_1(B)+\rk\pi_1(\tilde A)=\rk\pi_1(Y_N)\le2g$, so $|p|=N\le g$.
\end{proof}

\begin{remark}
The bound is attained: for $T(2,3)$ and $p=1$, $Y$ is the Poincar\'e sphere and $2=2$. For a nontrivial fibred knot $J$ and $p\ge2$, the cable $C_{p,1}(J)$ is fibred of genus $p\,g(J)\ge p$, and its $p$-surgery is $L(p,1)\#S^3_{1/p}(J)$ \cite{Gordon83}, in agreement with the bound. (This cable is not an L-space knot; Proposition~\ref{prop:cover} needs only fibredness.)
\end{remark}

\begin{remark}[avoiding Theorem~\ref{thm:perelman}]\label{rem:noperelman}
Replace rank by Heegaard genus $h$. Since $\beta\cdot\tilde\lambda=\pm1$, $Y_N$ is the closed manifold of an open book with page $\tilde F$, so $h(Y_N)\le2g$. By Haken's lemma \cite{Haken}, $h(Y_N)=h(\tilde A)+N h(B)$. Finally $h(B)\ge2$, since the closed $3$-manifolds of Heegaard genus at most $1$ are $S^3$, $S^1\times S^2$ and lens spaces. This again gives $N\le g$.
\end{remark}

\section{Heegaard Floer input}\label{sec:HF}

In this section $K$ is a nontrivial positive L-space knot of genus $g$, and $p\ne0$ is an integer with $N:=|p|$.

\subsection{Coset invariance}

Suppose $S^3_p(K)\cong L\#Y''$, where $L$ is an L-space with $a=|H_1(L)|\ge2$ and $Y''$ is a closed oriented $3$-manifold. If the homeomorphism reverses orientation, replace $(L,Y'')$ by $(-L,-Y'')$; $-L$ is again an L-space of the same order. Put $b:=N/a$. Then $b$ is an integer and $b\not\equiv0\pmod N$.

\begin{lemma}[coset invariance]\label{lem:coset}
For every label $i$, $\HFp(S^3_p(K),i)\cong\HFp(S^3_p(K),i+b)$ as relatively graded $\F[U]$-modules, up to a degree shift.
\end{lemma}

\begin{proof}
$H_1(L)$ is a subgroup of order $a$ of $\Z/N=\langle[\mu]\rangle$, so it is $\langle b[\mu]\rangle$. By \ref{H:sum}, the $\Spinc$ structures with a fixed $Y''$-component $\mathfrak t''$ form one $H_1(L)$-orbit, and all of them have $\HFp\cong\HFp(Y'',\mathfrak t'')$ up to shift. By \ref{H:affine}, the $H_1(L)$-orbits are exactly the cosets $i+b\Z$ of labels.
\end{proof}

\subsection{The invariant \texorpdfstring{$\tau$}{tau}}

\begin{definition}
Let $M$ be a graded $\F[U]$-module that is finite-dimensional in each degree. Record degrees by \emph{level}, meaning half the grading, so that $U$ lowers the level by $1$. Define $\tau(M)$ to be the multiset of levels $k$, where $k$ appears with multiplicity $\dim(M_k/U\cdot M_{k+1})$.
\end{definition}

$\tau$ is an isomorphism invariant, and shifting the grading translates $\tau$. If $M=\Tp\oplus\bigoplus_l\F[U]/U^{e_l}$, then $\tau(M)$ is the multiset of the top levels of the finite summands, because $U$ maps onto $\Tp$; in particular $|\tau(M)|$ is the number of finite summands. Clearly $\tau(M\oplus M')=\tau(M)\cup\tau(M')$.

\subsection{The level model and the structure of each class}

Fix a class $c\in\Z/N$. Let $s_1<s_2<\cdots<s_m$ be the integers $s\equiv c\pmod N$ with $|s|\le g-1$. Consecutive ones differ by $N$, and
\[
m=n_c:=\#\{s\equiv c\pmod N:\ |s|\le g-1\}.
\]

\paragraph{Level model.} By \ref{H:cone} we may compute with the homology-level maps $U^{V_s}$ and $U^{H_s}$ between copies of $\Tp$.
\begin{itemize}[leftmargin=*]
\item Record each tower $\Tp$ by its \emph{bottom level}.
\item Regrade each $B$-tower so that the cone maps preserve levels. This is possible: in one class the $A$'s and $B$'s form a connected graph (each $B$ is joined to its at most two adjacent $A$'s), and the relative grading of the cone is homogeneous. A map $U^e$ from a tower with bottom $x$ then lands in a tower with bottom $x+e$.
\item Under this regrading, an $A$-element at level $k$ maps to the $B$-element at level $k$ in each adjacent $B$ that is \emph{present at level $k$}, meaning that its bottom is $\le k$. All exponents are $\ge0$, so every present $B$ has its adjacent $A$'s present.
\item So at level $k$ the cone map is the coboundary $\delta_k\colon\F^{\text{vertices}}\to\F^{\text{edges}}$ of a graph $G_k$. Its vertices are the $A$'s with bottom $\le k$, and its edges are the $B$'s with bottom $\le k$, each $B$ joining its adjacent $A$'s.
\item $U$ acts on cochains by restriction from $G_{k+1}$ to $G_k$.
\item The cone has two terms and $D$ commutes with $U$, so $H(\mathrm{cone})=\ker D\oplus\operatorname{coker}D$ as graded $\F[U]$-modules.
\end{itemize}
The true relative grading is $2\cdot\text{level}+\text{constant}$ on each side, because both satisfy the same constraints on a connected graph. So $\tau$, computed in levels, is the invariant $\tau$ above up to translation.

\begin{proposition}[structure of a class]\label{prop:structure}
Let $c$ and $s_1<\dots<s_m$ be as above.
\begin{enumerate}[label=(\Alph*),leftmargin=*]
\item Let $0<p\le2g-2$ and $m\ge1$. Then $\tau(\HFp(S^3_p(K),c))=\{y-1:\ y\in Y_c\}$, where $Y_c$ is the multiset of the $m-1$ \emph{interior peak} levels
\[
y_j=\frac{\mu_j^2-p^2/4}{2p},\qquad \mu_j=s_j+\frac p2\qquad(1\le j\le m-1).
\]
In particular $|\tau|=m-1$.
\item Let $p=-N<0$ and $m\ge1$. Let $\Pi_c$ be the multiset of the $m+1$ \emph{peak} levels
\[
b(t)=-\frac{u^2-N^2/4}{2N},\qquad u=t+\frac N2,\qquad t\in\{s_1-N,\,s_1,\,s_2,\dots,s_m\},
\]
and let $M$ be its maximum. Then $\tau(\HFp(S^3_p(K),c))$ is $\{y-1:\ y\in\Pi_c\}$ with one copy of $M-1$ removed. In particular $|\tau|=m$.
\end{enumerate}
Neither multiset depends on $K$.
\end{proposition}

\begin{proof}[Proof of (A)]
In class $c$ the $A$'s and $B$'s form one infinite path $\cdots A_{s-p}-B_s-A_s-B_{s+p}-\cdots$. Writing $x(s)$ for the bottom of $A_s$, the bottoms must satisfy $x(s+p)-x(s)=H_s-V_{s+p}$, since $B_{s+p}$ receives $U^{H_s}$ from $A_s$ and $U^{V_{s+p}}$ from $A_{s+p}$. One solution is
\[
x(s)=\frac{s(s-p)}{2p}-V_s;
\]
indeed $x(s+p)-x(s)=s-V_{s+p}+V_s=H_s-V_{s+p}$, using $H_s=V_s+s$. The $B$ between $s$ and $s+p$ has bottom
\[
y=x(s)+H_s=\frac{s(s+p)}{2p}=\frac{\mu^2-p^2/4}{2p},\qquad \mu=s+\frac p2.
\]
We have $x(s)\to+\infty$ as $|s|\to\infty$: for $s\ge g$, $V_s=0$, and for $s<0$, $V_s=V_{-s}+|s|\le g+|s|$. So the quadratic term dominates, and each $G_k$ is a finite subforest of a path.

The coboundary of a finite forest is onto. So $D$ is surjective and $\HFp\cong\ker D$. At level $k$, $\ker\delta_k$ is the space of functions on the components of $G_k$, and $U$ is restriction of functions. Hence $\dim(M_k/U\cdot M_{k+1})$ is the number of \emph{merges}: the number of components of $G_k$ minus the number of components of $G_{k+1}$ that contain a vertex of $G_k$.

\emph{Edges with a zero exponent.} $V_s=0$ iff $s\ge g$, and $H_s=0$ iff $s\le -g$. A $B$ with both exponents $0$ would be $B_{s+p}$ with $s\le-g$ and $s+p\ge g$, so $p\ge2g$, which is excluded. A $B$ with exactly one exponent $0$ appears at the same level as its endpoint $A_e$ on that side. That endpoint is new at that level. Its other edge has a positive exponent at $A_e$, since $V_e=H_e=0$ would need $e\ge g$ and $e\le-g$; so that other edge appears strictly later. The far endpoint of $B$ is old, since $B$'s exponent there is positive. So at this level $B$ attaches the new vertex $A_e$ to one old component, and merges no old components.

\emph{Edges with both exponents $\ge1$.} The edge $B_{s+p}$ between $A_s$ and $A_{s+p}$ has exponents $H_s\ge1$ (iff $s\ge-(g-1)$) and $V_{s+p}\ge1$ (iff $s+p\le g-1$). So these are exactly the $m-1$ edges between consecutive $s_j,s_{j+1}$: the interior peaks. Each appears strictly after both of its endpoints, and since $G$ is a path, each one joins two distinct components of the old graph.

So the merges from level $k$ to level $k+1$ are the interior peaks at level $k+1$. Hence $\tau=\{y-1:\ y\text{ an interior peak}\}$ and $|\tau|=m-1$.
\end{proof}

\begin{proof}[Proof of (B)]
Here $h_s\colon A_s\to B_{s-N}$. Let $C$ be the class-$c$ cone.

\emph{Truncation, step 1.} Let $P=\{A_s:s<g\}\cup\{B_t:t<g\}$. It is a subcomplex: for $s<g$, $A_s$ maps to $B_s$ and $B_{s-N}$, and both indices are $<g$. The quotient $Q=C/P=\{A_s,B_s:s\ge g\}$ has differential $D_Q=v_s+h_s$, where $v_s\colon A_s\to B_s$ is the identity since $V_s=0$, and $h_s$ lowers the index by $N$ and is present only while $s-N\ge g$. On finitely supported chains, the identity plus a part that lowers the index, with indices bounded below, is bijective: invert it by a finite Neumann series. So $Q$ is acyclic and $H(C)\cong H(P)$.

\emph{Truncation, step 2.} Let $P'=\{A_s:|s|<g\}\cup\{B_t:-g-N<t<g\}$. It is a subcomplex of $P$: for $|s|<g$ we have $s-N>-g-N$. The quotient $R_0=P/P'=\{A_s:s\le-g\}\cup\{B_t:t\le-g-N\}$ has differential $h_s+v_s$, where $h_s\colon A_s\to B_{s-N}$ is the identity since $H_s=0$, and $v_s\colon A_s\to B_s$ is present only while $s\le-g-N$. Relative to the pairing $A_s\leftrightarrow B_{s-N}$, the second part raises the $B$-index by $N$, and indices are bounded above. So it is again identity plus nilpotent on finitely supported chains, and $R_0$ is acyclic. Hence $H(C)\cong H(P')$.

\emph{Shape of $P'$.} The class-$c$ integers in $(-g-N,g)$ are $s_1-N,s_1,\dots,s_m$, because $s_1-N\le-g<s_1$ and $s_m<g\le s_m+N$. So $P'$ is the zigzag
\[
B_{s_1-N}-A_{s_1}-B_{s_1}-A_{s_2}-\cdots-A_{s_m}-B_{s_m}.
\]
Every exponent on it is $\ge1$: $V_{s_j}\ge1$ since $s_j\le g-1$, and $H_{s_j}=V_{-s_j}\ge1$ since $-s_j\le g-1$. The two end $B$'s each have one $A$-neighbour. Treat them as edges to a \emph{ground} vertex, at which cochains take the value $0$. Then $G\cup\{\text{ground}\}$ is a single cycle with $m+1$ edges.

\emph{Levels.} Put $x(s)=-\frac{s(s+N)}{2N}-V_s$ and $b(t)=-\frac{t(t+N)}{2N}=-\frac{u^2-N^2/4}{2N}$, $u=t+\frac N2$. Then $B_s$ has bottom $x(s)+V_s=b(s)$, and since $H_s=V_s+s$,
\[
x(s)+H_s=b(s)+s=\frac{-s^2-sN+2sN}{2N}=-\frac{s(s-N)}{2N}=b(s-N),
\]
as required for $B_{s-N}$.

\emph{The cokernel part.} $\operatorname{coker}\delta_k=H^1(G_k\cup\text{ground},\text{ground})$, which is $H^1(G_k\cup\text{ground})$ because the ground is one point. This is $0$ until the cycle closes, at the maximal peak level $M$. From then on it is $\F$, and $U$ (restriction of cochains) is an isomorphism. So the cokernel part is a single tower and contributes nothing to $\tau$.

\emph{The kernel part.} $\ker\delta_k$ is the space of functions on the components of $G_k$ that do not contain the ground. A vertex that is new at level $k+1$ has its edges strictly higher, since all exponents are $\ge1$. So it is isolated at level $k+1$, and the new edges at level $k+1$ join only old vertices or the ground. Hence $\dim(\ker_k/U\cdot\ker_{k+1})$ is the number of merges among old components, including merges into the ground component. That number is $\#(\text{new edges at level }k+1)-\Delta b_1$, and $b_1$ increases (by $1$) only at level $M$. So $\tau=\{y-1:y\in\Pi_c\}$ with one copy of $M-1$ removed, and $|\tau|=(m+1)-1=m$.
\end{proof}

\subsection{Divisibility}

\begin{proposition}\label{prop:divis}
Suppose $2\le|p|\le2g-2$ and $S^3_p(K)\cong L\#Y''$ with $L$ an L-space of order $a\ge2$. Then $|p|$ divides $2g-1$.
\end{proposition}

\begin{proof}
Write $2g-1=QN+r$ with $0\le r<N$. Then $Q\ge1$, since $N\le2g-2$. The integers in $[-(g-1),g-1]$ are $2g-1$ consecutive integers, so $n_i=Q+1$ on a set $S$ of $r$ cyclically consecutive residues, and $n_i=Q$ elsewhere. In particular every $n_i\ge1$.

By Proposition~\ref{prop:structure}, $|\tau(\HFp(S^3_p(K),i))|$ equals $n_i-1$ if $p>0$ and $n_i$ if $p<0$; either way it is a strictly increasing function of $n_i$. By Lemma~\ref{lem:coset} and the invariance of $\tau$ up to translation, $|\tau|$ is invariant under $i\mapsto i+b$ with $b\not\equiv0$. So $S$ is invariant under translation by $b$.

If $0<r<N$, then $S$ is a proper nonempty cyclic interval in $\Z/N$. It has exactly one start point (an $x\in S$ with $x-1\notin S$), and a translation preserving $S$ must fix that start point. That is impossible for $b\not\equiv0$. So $r=0$, i.e.\ $N$ divides $2g-1$.
\end{proof}

\subsection{The key contradiction}

\begin{theorem}\label{thm:key}
Let $K$ be a nontrivial positive L-space knot of genus $g$, and $p$ a nonzero integer with $|p|\ge2$, $|p|$ dividing $2g-1$ and $|p|<2g-1$. Then $S^3_p(K)$ is not homeomorphic to $L\#Y''$ for any L-space $L$ with $|H_1(L)|\ge2$ and any closed oriented $3$-manifold $Y''$.
\end{theorem}

\begin{proof}
Let $N:=|p|$ and $n:=(2g-1)/N$. Both are odd, $N\ge3$ and $n\ge3$. In particular $|p|\le(2g-1)/3\le2g-2$, so Proposition~\ref{prop:structure} applies. Suppose such a splitting exists, and let $b$ be as in Lemma~\ref{lem:coset}.

\emph{(i) The comparison.} Let $t$ be the representative of $b$ with $0<|t|\le(N-1)/2$; it exists because $N$ is odd and $b\not\equiv0$. By Lemma~\ref{lem:coset} and the invariance of $\tau$:
\[
(\star)\qquad \tau\text{ of class }0\text{ and }\tau\text{ of class }t\text{ agree up to a translation.}
\]

\emph{(ii) Arithmetic.} For $|c|\le(N-1)/2$, the integers $s\equiv c$ with $|s|\le g-1$ are exactly $c+kN$ for $|k|\le(n-1)/2$. Indeed $|c+kN|\le(N-1)/2+(n-1)N/2=g-1$, while $|k|\ge(n+1)/2$ gives $|c+kN|\ge(nN+1)/2>g-1$. So $m=n$ in every class.

\emph{Key fact.} Let $u=c+(k+\tfrac12)N$ and $u'=c+(k'+\tfrac12)N$ with $k\ne k'$. If $u^2=u'^2$, then $u=-u'$, so $2c+(k+k'+1)N=0$. Hence $N\mid 2c$, so $N\mid c$ because $N$ is odd, and $c=0$. So for $c=t$, distinct $k$ give distinct values of $u^2$, hence distinct levels in both formulas of Proposition~\ref{prop:structure}.

\emph{(iii) Case $p>0$.} \emph{Class $0$:} $s_j=(j-\tfrac{n+1}2)N$, so the interior peaks have $\mu_j=s_j+N/2=(j-\tfrac n2)N$ for $1\le j\le n-1$. These are $\pm N/2,\pm3N/2,\dots,\pm(n-2)N/2$: nonzero $\pm$ pairs. So every element of $\tau$ has even multiplicity, and $\tau$ is nonempty since $n-1\ge2$. \emph{Class $t$:} $\mu_j=t+(j-\tfrac{n+1}2+\tfrac12)N$, and by (ii) the $n-1$ levels are distinct. A translate of a nonempty multiset in which every multiplicity is even never has all multiplicities equal to $1$. This contradicts $(\star)$.

\emph{(iv) Case $p<0$.} \emph{Class $0$:} the peaks are $t'=kN$ for $k=-\tfrac{n+1}2,\dots,\tfrac{n-1}2$, so $u=t'+N/2$ runs over $\pm N/2,\pm3N/2,\dots,\pm nN/2$, each value of $|u|$ exactly twice. The maximal level $M$ comes from $|u|=N/2$. Removing one copy of $M-1$ leaves the levels for $|u|=3N/2,\dots,nN/2$ still doubled, and these exist since $n\ge3$. So $\tau$ has a repeated element. \emph{Class $t$:} by (ii) the $n+1$ peak levels are distinct, so the $n$ elements of $\tau$ are distinct. This contradicts $(\star)$.
\end{proof}

\section{Proof of the Main Theorem}\label{sec:proof}

Assume, as in Section~\ref{sec:reduction}, that $K$ is a nontrivial positive L-space knot of genus $g$, $p\ne0$, and $Y=S^3_p(K)$ has $k\ge3$ prime summands.
\begin{enumerate}[leftmargin=*]
\item By Lemma~\ref{lem:hs}, $Y\cong A\#B$ with $B$ an integer homology sphere, $B\ne S^3$.
\item By Lemma~\ref{lem:lens}, $Y\cong L\#Y''$ with $L$ a lens space, hence an L-space, of order $a\ge2$. So $|p|\ge2$.
\item $K$ is fibred with fibre of genus $g$ and one boundary component (Theorem~\ref{thm:fibred}). Proposition~\ref{prop:cover} gives $|p|\le g$. If $g=1$, this contradicts $|p|\ge2$. So $g\ge2$, and $2\le|p|\le g\le2g-2$.
\item By Proposition~\ref{prop:divis}, applied to the splitting of step~2, $|p|$ divides $2g-1$. Also $|p|\le g<2g-1$.
\item By Theorem~\ref{thm:key}, the splitting of step~2 cannot exist. This is a contradiction.
\end{enumerate}
So $S^3_r(K)$ has at most two prime summands, for every L-space knot $K$ and every slope $r$.\qed

\section{Remarks}

\begin{itemize}[leftmargin=*]
\item For a fibred knot, a reducible integral surgery $S^3_p(K)$ with an integer homology sphere summand $B\ne S^3$ requires $|p|\le g$ (Proposition~\ref{prop:cover}). This uses only fibredness.
\item Granting Theorem~\ref{thm:GL96}, the argument of Section~\ref{sec:proof} shows more: a \emph{reducible} surgery on an L-space knot never has a prime summand that is an integer homology sphere other than $S^3$, since steps 2--5 use only the homology sphere summand and the lens space summand. (Irreducible surgeries can be homology spheres, e.g.\ $S^3_{+1}(T(2,3))$ is the Poincar\'e sphere.) This agrees with the known examples, cables at the cabling slope, whose summands are $L(a,b)$ and $S^3_{b/a}(J)$.
\end{itemize}

\begin{thebibliography}{99}

\bibitem{Gabai}
D.~Gabai, Foliations and the topology of $3$-manifolds.~III, \emph{J. Differential Geom.} \textbf{26} (1987), 479--536.

\bibitem{Ghiggini}
P.~Ghiggini, Knot Floer homology detects genus-one fibred knots, \emph{Amer. J. Math.} \textbf{130} (2008).

\bibitem{Gordon83}
C.~McA.~Gordon, Dehn surgery and satellite knots, \emph{Trans. Amer. Math. Soc.} \textbf{275} (1983), 687--708.

\bibitem{GL87}
C.~McA.~Gordon and J.~Luecke, Only integral Dehn surgeries can yield reducible manifolds, \emph{Math. Proc. Cambridge Philos. Soc.} \textbf{102} (1987), 97--101.

\bibitem{GL96}
C.~McA.~Gordon and J.~Luecke, Reducible manifolds and Dehn surgery, \emph{Topology} \textbf{35} (1996), 385--409.

\bibitem{Grushko}
I.~A.~Grushko, On the bases of a free product of groups, \emph{Mat. Sbornik} \textbf{8} (1940).

\bibitem{Haken}
W.~Haken, Some results on surfaces in $3$-manifolds, in: \emph{Studies in Modern Topology}, MAA Studies in Mathematics~5, 1968.

\bibitem{Howie}
J.~Howie, A proof of the Scott--Wiegold conjecture on free products of cyclic groups, \emph{J. Pure Appl. Algebra} \textbf{173} (2002).

\bibitem{Kneser}
H.~Kneser, Geschlossene Fl\"achen in dreidimensionalen Mannigfaltigkeiten, \emph{Jahresber. Deutsch. Math.-Verein.} \textbf{38} (1929).

\bibitem{LS}
R.~C.~Lyndon and P.~E.~Schupp, \emph{Combinatorial Group Theory}, Springer, 1977.

\bibitem{Milnor}
J.~Milnor, A unique decomposition theorem for $3$-manifolds, \emph{Amer. J. Math.} \textbf{84} (1962), 1--7.

\bibitem{MorganTian}
J.~Morgan and G.~Tian, \emph{Ricci Flow and the Poincar\'e Conjecture}, Clay Mathematics Monographs~3, Amer. Math. Soc., 2007.

\bibitem{Moser}
L.~Moser, Elementary surgery along a torus knot, \emph{Pacific J. Math.} \textbf{38} (1971), 737--745.

\bibitem{Neumann}
B.~H.~Neumann, On the number of generators of a free product, \emph{J. London Math. Soc.} \textbf{18} (1943).

\bibitem{Ni}
Y.~Ni, Knot Floer homology detects fibred knots, \emph{Invent. Math.} \textbf{170} (2007), 577--608.

\bibitem{NiWu}
Y.~Ni and Z.~Wu, Cosmetic surgeries on knots in $S^3$, \emph{J. Reine Angew. Math.} \textbf{706} (2015).

\bibitem{OSapp}
P.~Ozsv\'ath and Z.~Szab\'o, Holomorphic disks and three-manifold invariants: properties and applications, \emph{Ann. of Math. (2)} \textbf{159} (2004), 1159--1245.

\bibitem{OSlens}
P.~Ozsv\'ath and Z.~Szab\'o, On knot Floer homology and lens space surgeries, \emph{Topology} \textbf{44} (2005), 1281--1300.

\bibitem{OSinteger}
P.~Ozsv\'ath and Z.~Szab\'o, Knot Floer homology and integer surgeries, \emph{Algebr. Geom. Topol.} \textbf{8} (2008), 101--153.

\bibitem{Perelman1}
G.~Perelman, The entropy formula for the Ricci flow and its geometric applications, arXiv:math/0211159, 2002.

\bibitem{Perelman2}
G.~Perelman, Ricci flow with surgery on three-manifolds, arXiv:math/0303109, 2003.

\bibitem{Perelman3}
G.~Perelman, Finite extinction time for the solutions to the Ricci flow on certain three-manifolds, arXiv:math/0307245, 2003.

\end{thebibliography}

\end{document}
